%! Confidence Intervals for the Difference of Two Proportions — Unit 6, Lesson 6.8 — Group Activity
\documentclass[10pt]{article}
\usepackage{apstats-article}
\usepackage{apstats-boxes}
\newcommand{\TallMath}[1]{$\displaystyle #1\rule[-1.4em]{0pt}{3.2em}$}

\begin{document}

\pageheader{Unit 6, Lesson 6.8}{Group Activity}
\namepartnerperiod

\vspace{0.15in}

\begin{scenariobox}[Scenario 1 --- Social Media Use by Grade Level]{sky}
A researcher surveyed random samples of high school students to compare social media use.
Among 180 sophomores, 126 reported using social media more than 3 hours per day.
Among 210 seniors, 168 reported the same.

\begin{enumerate}[label=\arabic*.]
  \item Identify the two populations and define $p_1$ and $p_2$.
        \writeline

  \item Compute the point estimate $\hat p_1 - \hat p_2$.
        \writeline

  \item Verify all three conditions (Random, Large Counts, 10\%).
        \writelines{4}

  \item Compute $SE_{\hat p_1 - \hat p_2}$.
        \writelines{2}

  \item Construct a 95\% confidence interval.
        \writelines{2}

  \item Interpret the interval in context.
        \writelines{2}
\end{enumerate}
\end{scenariobox}

\begin{scenariobox}[Scenario 2 --- Flu Vaccination Rates]{sky}
A public health department took independent random samples to compare flu vaccination rates
between two counties. In County A ($n_1 = 300$), 204 residents had received a flu vaccine.
In County B ($n_2 = 250$), 145 had.

\begin{enumerate}[label=\arabic*.]
  \item Compute $\hat p_A$ and $\hat p_B$, then the point estimate.
        \writeline

  \item Verify the Large Counts condition (show all four values).
        \writelines{2}

  \item Construct a 99\% confidence interval for $p_A - p_B$.
        \writelines{3}

  \item Interpret the interval, then answer: does this interval provide evidence that County A's
        vaccination rate is higher? Explain.
        \writelines{2}
\end{enumerate}
\end{scenariobox}

\end{document}
